\documentclass[UTF8,zihao=-4]{ctexart}
\usepackage[a4paper,margin=2.1cm]{geometry}
\usepackage{hyperref}
\usepackage{array}
\usepackage{longtable}
\usepackage{verbatim}
\hypersetup{hidelinks}
\setlength{\parindent}{0pt}
\setlength{\parskip}{0.55em}
\emergencystretch=4em
\sloppy
\title{02-语法分析复习讲义}
\author{}
\date{}
\begin{document}
\maketitle
\setcounter{tocdepth}{1}
\tableofcontents
\newpage
\section{02 语法分析复习讲义}

对应资料：

\begin{itemize}
\item \texttt{Slides/Slides/Ch4-1-Context-Free Language and Push-down Automata.pdf}
\item \texttt{Slides/Slides/Ch4-2-Syntax Analysis.pdf}
\item \texttt{Slides/Slides/Slides-Ant/4-parser-cfg-antlr-handout.pdf}
\item \texttt{Slides/Slides/Slides-Ant/5-parser-ll-handout.pdf}
\item \texttt{Slides/Slides/Slides-Ant/6-parser-allstar-handout.pdf}
\item \texttt{Slides/Slides/Slides-Ant/12-parser-lr0.pdf}
\item \texttt{Review/compile.pdf}
\item \texttt{Exam/2025.pdf} 语法分析题
\item \texttt{Exam/过时/2022-exam-A.pdf} 题目 2、3、4
\item \texttt{Exam/过时/2023-exam-A.pdf} 题目 2、3
\end{itemize}

\subsection{一、CFG}

上下文无关文法：

\begin{quote}
\footnotesize
\begin{verbatim}
G = (N, T, S, P)
N : 非终结符集合
T : 终结符集合，N cap T = emptyset
S : 起始符号，S in N
P : 产生式集合
\end{verbatim}
\end{quote}

派生：

\begin{quote}
\footnotesize
\begin{verbatim}
xAy => x alpha y       使用 A -> alpha 展开
=>*                    多步派生
=>lm                   最左派生
=>rm                   最右派生
\end{verbatim}
\end{quote}

语言：

\begin{quote}
\footnotesize
\begin{verbatim}
L(G) = {w in T* | S =>* w}
\end{verbatim}
\end{quote}

\subsection{二、推导树与二义性}

推导树记录产生式展开结构。同一个终结符串有两棵不同推导树时，文法有二义性。

典型二义性：

\begin{quote}
\footnotesize
\begin{verbatim}
E -> E + E | E * E | a
\end{verbatim}
\end{quote}

消除优先级和结合性后的写法：

\begin{quote}
\footnotesize
\begin{verbatim}
E -> E + T | T
T -> T * F | F
F -> a | (E)
\end{verbatim}
\end{quote}

\subsection{三、CFL 与 NPDA}

NPDA 定义：

\begin{quote}
\footnotesize
\begin{verbatim}
P = (Q, Sigma, Gamma, delta, q0, z0, F)
Q      : 状态集合
Sigma  : 输入字母表
Gamma  : 栈字母表
q0     : 起始状态
z0     : 起始栈符号
F      : 终止状态集合
\end{verbatim}
\end{quote}

结论：

\begin{itemize}
\item CFG 与 NPDA 定义同一类语言，即上下文无关语言。
\item CFL 对并、连接、反转、Kleene star 封闭。
\item CFL 对交、补、差不封闭。
\item NPDA 比 DPDA 表达能力更强。
\end{itemize}

\subsection{四、CFL 泵引理}

形式：

\begin{quote}
\footnotesize
\begin{verbatim}
若 L 是 CFL，则存在 C，
对任意 s in L 且 |s| >= C，
存在 s = uvwxy，满足 vx != epsilon，|vwx| <= C，
对任意 i >= 0，u v^i w x^i y in L。
\end{verbatim}
\end{quote}

用途：证明语言不是 CFL。

\subsection{五、左递归消除}

直接左递归：

\begin{quote}
\footnotesize
\begin{verbatim}
A -> A alpha1 | A alpha2 | beta1 | beta2
\end{verbatim}
\end{quote}

改写为：

\begin{quote}
\footnotesize
\begin{verbatim}
A  -> beta1 A' | beta2 A'
A' -> alpha1 A' | alpha2 A' | epsilon
\end{verbatim}
\end{quote}

例：

\begin{quote}
\footnotesize
\begin{verbatim}
E -> E a | b
\end{verbatim}
\end{quote}

消除后：

\begin{quote}
\footnotesize
\begin{verbatim}
E  -> b E'
E' -> a E' | epsilon
\end{verbatim}
\end{quote}

考试要点：

\begin{itemize}
\item LL 分析前先消除左递归。
\item 左递归消除后，FIRST/FOLLOW 和预测表要基于新文法计算。
\end{itemize}

\subsection{六、FIRST 与 FOLLOW}

FIRST：

\begin{quote}
\footnotesize
\begin{verbatim}
FIRST(alpha) = alpha 能推出的串的第一个终结符集合；
若 alpha =>* epsilon，则 epsilon in FIRST(alpha)。
\end{verbatim}
\end{quote}

FOLLOW：

\begin{quote}
\footnotesize
\begin{verbatim}
FOLLOW(A) = 在某个句型中紧跟 A 后面能出现的终结符集合；
起始符号 S 的 FOLLOW(S) 包含 $。
\end{verbatim}
\end{quote}

计算规则：

\begin{itemize}
\item 若 \texttt{A -> alpha B beta}，则 \texttt{FIRST(beta) - \{epsilon\}} 加入 \texttt{FOLLOW(B)}。
\item 若 \texttt{A -> alpha B}，或 \texttt{A -> alpha B beta} 且 \texttt{epsilon in FIRST(beta)}，则 \texttt{FOLLOW(A)} 加入 \texttt{FOLLOW(B)}。
\item 反复迭代直到集合不再变化。
\end{itemize}

\subsection{七、LL(1) 预测分析表}

对每条产生式 \texttt{A -> alpha}：

\begin{enumerate}
\item 对每个 \texttt{a in FIRST(alpha) - \{epsilon\}}，填：
\end{enumerate}

\begin{quote}
\footnotesize
\begin{verbatim}
M[A, a] = A -> alpha
\end{verbatim}
\end{quote}

\begin{enumerate}
\item 若 \texttt{epsilon in FIRST(alpha)}，对每个 \texttt{b in FOLLOW(A)}，填：
\end{enumerate}

\begin{quote}
\footnotesize
\begin{verbatim}
M[A, b] = A -> alpha
\end{verbatim}
\end{quote}

LL(1) 判断：

\begin{itemize}
\item 表中每个格子最多一条产生式，就是 LL(1)。
\item 若某格多条产生式，存在冲突，不是 LL(1)。
\end{itemize}

\subsection{八、LR 分析核心}

自底向上分析基于栈：

\begin{itemize}
\item shift：移入输入符号。
\item reduce：把栈顶一段右部归约为左部。
\item accept：归约到起始符号并读到 \texttt{\$}。
\item error：无合法动作。
\end{itemize}

LR(0) 项：

\begin{quote}
\footnotesize
\begin{verbatim}
A -> alpha . beta
\end{verbatim}
\end{quote}

closure：

\begin{quote}
\footnotesize
\begin{verbatim}
若 A -> alpha . B beta 在 I 中，
对 B -> gamma，把 B -> . gamma 加入 I。
\end{verbatim}
\end{quote}

goto：

\begin{quote}
\footnotesize
\begin{verbatim}
goto(I, X) = closure({A -> alpha X . beta | A -> alpha . X beta in I})
\end{verbatim}
\end{quote}

\subsection{九、SLR 与 LR(1)}

SLR 表：

\begin{itemize}
\item 对 \texttt{goto(I, a)=J}，终结符 \texttt{a} 填 shift \texttt{sJ}。
\item 对 \texttt{A -> alpha .}，在 \texttt{FOLLOW(A)} 的每个终结符上填 reduce \texttt{A -> alpha}。
\item 对 \texttt{S' -> S .}，在 \texttt{\$} 上填 accept。
\item 若同一格有多个动作，出现冲突。
\end{itemize}

LR(1) 项：

\begin{quote}
\footnotesize
\begin{verbatim}
[A -> alpha . beta, a]
\end{verbatim}
\end{quote}

其中 \texttt{a} 是展望符。LR(1) 归约只在展望符对应列填 reduce，精度高于 SLR。

\subsection{十、ANTLR4 与优先级上升}

ANTLR4 可以直接写左递归表达式规则，然后改写成带优先级参数的规则。

常见模式：

\begin{quote}
\footnotesize
\begin{verbatim}
expr
  : expr postfix
  | expr '^' expr
  | expr '+' expr
  | INT
\end{verbatim}
\end{quote}

答题要点：

\begin{itemize}
\item 后缀运算符优先级最高。
\item 右结合运算符在递归调用中使用同级或更低约束。
\item 左结合运算符右侧递归调用使用更高优先级约束。
\item 画树时先按优先级分组，再按结合性确定嵌套方向。
\end{itemize}

\subsection{十一、自测题与答案}

\subsubsection{题 1：\texttt{A -> A a | b} 消除左递归。}

答案：

\begin{quote}
\footnotesize
\begin{verbatim}
A  -> b A'
A' -> a A' | epsilon
\end{verbatim}
\end{quote}

\subsubsection{题 2：LL(1) 表冲突意味着什么？}

答案：同一个非终结符和同一个展望符下有多条可用产生式，分析器无法只看 1 个 token 决定展开规则。

\subsubsection{题 3：SLR 与 LR(1) 的主要区别是什么？}

答案：SLR 用 FOLLOW 集合设置归约项；LR(1) 在每个项目中携带展望符，只在项目展望符上设置归约，冲突更少。

\end{document}
